% from ega4-errata-addenda-list3-fr.tex
\oldpage[IV]{351}
plat par ($0_{\mathrm I}$, 6.1.2); enfin, les $H_i(L_\bullet)$ sont nuls
par définition sauf $H_0(L_\bullet)=N$, qui est un $A$-module plat. On
sait que, dans ces conditions,
\[
 H_n(K_\bullet\otimes_AL_\bullet)\simeq
 \bigoplus_{p+q=n}
 H_p(K_\bullet)\otimes_AH_q(L_\bullet)
\]
(G, I, 5.5), donc $K_\bullet\otimes_AL_\bullet$ (gradué comme
d'ordinaire) est une résolution de $M\otimes_AN$ par des $B$-modules
libres. Posons
\[
 r=\operatorname{dim.proj}_A(M),\qquad
 s=\operatorname{dim.proj}_B(N)=
 \operatorname{dim.proj}_{B\otimes_Ak}(N\otimes_Ak).
\]
Il y a donc une résolution libre $K_\bullet$ de longueur $r$ et une
résolution libre $L_\bullet$ de longueur $s$, et
$K_\bullet\otimes_AL_\bullet$ est alors une résolution libre de longueur
$r+s$, ce qui prouve déjà que
$\operatorname{dim.proj}_B(M\otimes_AN)\leq r+s$. Il reste à prouver
($0$, 17.2.6) que l'on a
\[
 \operatorname{Tor}^{B}_{r+s}(M\otimes_AN,k(B))\neq0.
\]
Utilisant la résolution précédente $K_\bullet\otimes_AL_\bullet$ pour
calculer ce foncteur, on obtient
\[
 \operatorname{Tor}^{B}_{r+s}(M\otimes_AN,k(B))=
 H_{r+s}((K_\bullet\otimes_AL_\bullet)\otimes_Bk(B)).
\]
Mais on peut écrire
\[
 (K_\bullet\otimes_AL_\bullet)\otimes_Bk(B)=
 (K_\bullet\otimes_Ak)\otimes_k
 (L_\bullet\otimes_Bk(B)),
\]
qui est donc un produit tensoriel de deux complexes d'espaces
vectoriels sur $k$; tous les $k$-modules étant plats, la formule de
Künneth donne
\[
 H_{r+s}((K_\bullet\otimes_Ak)\otimes_k
 (L_\bullet\otimes_Bk(B)))\simeq
 H_r(K_\bullet\otimes_Ak)\otimes_k
 H_s(L_\bullet\otimes_Bk(B)),
\]
l'homologie du premier (resp. second) complexe étant nulle en degrés
$>r$ (resp. $>s$). On a donc
\[
 \operatorname{Tor}^{B}_{r+s}(M\otimes_AN,k(B))=
 \operatorname{Tor}^{A}_{r}(M,k)\otimes_k
 \operatorname{Tor}^{B}_{s}(N,k(B))
\]
à isomorphie près, et comme chacun des facteurs est $\neq0$ par
hypothèse, il en est de même de leur produit tensoriel sur $k$.

\paragraph{(Err\textsubscript{IV}, 23).}
Dans (IV, 6.10.1), ligne 8 de la p.~155, après : canonique, ajouter :
$\sh{F}$ un $\sh{O}_X$-Module cohérent.

\paragraph{(Err\textsubscript{IV}, 24).}
Dans ({\hyperref[canon-index.ega4.c04.prop.6-14-4]{IV, 6.14.4}}), la référence
donnée ligne 12 du bas de la p.~175 est insuffisante; il faut raisonner
de la façon suivante. Tout d'abord, si $M$ est une base de transcendance
de $K'$ sur $A$, $A(M)$ est algébriquement fermé dans $B(A(M))$
(Bourbaki, \emph{Alg.}, chap.~V, \S6, exerc.~8 c)). On peut donc se
borner au cas où $K'$ est une extension algébrique séparable de $A$.
Soit $p$ l'exposant caractéristique de $A$; il résulte de Bourbaki,
\emph{loc. cit.}, \S9, exerc.~2, que la fermeture algébrique de $K'$
dans $B(K')$ est radicielle sur $K'$; montrons que tout élément
$x\in B(K')$ tel que $x^{p^r}\in K'$ appartient déjà à $K'$. Pour cela,
considérons une base $(e_\lambda)$ de $K'$ sur $A$; puisque $K'$ est une
extension algébrique séparable de $A$, les $e_\lambda^{p^r}$ forment
aussi une base de $K'$ sur $A$ (Bourbaki, \emph{loc. cit.}, \S8,
n\textsuperscript{o}~3, cor. de la prop.~4). On a
$x=\sum_\lambda e_\lambda b_\lambda$, avec $b_\lambda\in B$ et, par
hypothèse, il existe des éléments $a_\lambda\in A$ tels que l'on ait
\[
 \sum_\lambda e_\lambda^{p^r}b_\lambda^{p^r}=
 \sum_\lambda e_\lambda^{p^r}a_\lambda;
\]
mais cela entraîne $b_\lambda^{p^r}=a_\lambda$ pour tout $\lambda$;
comme par hypothèse $A$ est algébriquement fermé dans $B$, on a donc
$b_\lambda\in A$ pour tout $\lambda$, et par suite $x\in K'$.

\paragraph{(Err\textsubscript{IV}, 25).}
Dans (IV, 7.4.8, B)), p.~203, la question a été résolue affirmativement
par R.~Kiehl, \emph{Ausgezeichnete Ringe der nichtarchimedischen
analytischen Geometrie}, à paraître dans \emph{Inventiones
Mathematicae}.

\paragraph{(Err\textsubscript{IV}, 26).}
Dans (IV, 7.8.4), ligne 9 du bas de la p.~217, remplacer :
\begin{quote}
Nous verrons plus loin (\S18) que si $k$ est un corps valué complet non
discret\ldots
\end{quote}
par :
\begin{quote}
Il a été démontré par R.~Kiehl (dans l'article cité en
(Err\textsubscript{IV}, 25)) que si $k$ est un corps valué complet non
discret\ldots
\end{quote}

